\documentclass{article}%
\usepackage[T1]{fontenc}%
\usepackage[utf8]{inputenc}%
\usepackage{lmodern}%
\usepackage{textcomp}%
\usepackage{lastpage}%
\usepackage{geometry}%
\geometry{margin=1in,headheight=14pt,headsep=25pt}%
\usepackage{hyperref}%
\usepackage{enumitem}%
\usepackage{fancyhdr}%
\usepackage{xcolor}%
\usepackage{url}%
\usepackage{breakurl}%
\usepackage{amsmath}%
\usepackage{amssymb}%
\usepackage{mathtools}%
\usepackage{amsthm}%
\usepackage{thmtools}%
\usepackage{float}%
\usepackage{subcaption}%
\usepackage[font=small,labelfont=bf]{caption}%
%

            \hypersetup{
                colorlinks=true,
                linkcolor=blue,
                filecolor=magenta,
                urlcolor=blue,
                breaklinks=true,
                bookmarks=true,
                pdfborder={0 0 0}
            }
            \urlstyle{same}
            \Urlmuskip=0mu plus 1mu  % URL breaking in nicer spots

            % Math configuration
            \everymath{\displaystyle}
            \setlength{\jot}{10pt}  % Increase spacing between equations

            \pagestyle{fancy}
            \fancyhf{}
            \rhead{Generated on \today}
            \cfoot{\thepage}

            % Better Unicode support
            \usepackage[utf8]{inputenc}
            \usepackage{newunicodechar}
            \setlength{\itemsep}{1em} % Adds spacing between items

            % Configure itemize settings globally
            \setlist[itemize]{nosep}
            \setlist[itemize]{leftmargin=*}

            % Remove section numbering
            \renewcommand{\thesection}{}
            \renewcommand{\thesubsection}{}

            % Adjust caption spacing around figures
            \setlength{\abovecaptionskip}{0pt}
            \setlength{\belowcaptionskip}{5pt}
        %
\lhead{Lecture: 2024{-}10{-}28{-}Farkas}%
\title{2024{-}10{-}28{-}Farkas}%
\author{Generated by Scribe.AI}%
\date{\today}%
%
\begin{document}%
\normalsize%
\maketitle%
\section{Slide 1}%
\label{sec:Slide1}%
\subsection{Content}%

\text{Farkas Lemma} \\
\text{LEMMA 10.5. The system } Ax \le b \text{ has no solutions if and only if there is a } y \text{ such} \\
\text{that} \\
(10.8) \quad A^T y = 0 \\
\quad y \ge 0 \\
\quad b^T y < 0 \\
\text{PROOF. Consider the linear program} \\
(P) \quad \text{maximize } 0 \\
\quad \text{subject to } Ax \le b \\
\text{and its dual} \\
(D) \quad \text{minimize } b^T y \\
\quad \text{subject to } A^T y = 0 \\
\quad y \ge 0
%
\subsection{Description}%
Slide 1 of 11 introduces Farkas' Lemma, a fundamental result in linear programming. The lemma states that a system of linear inequalities $Ax \le b$ has no solution if and only if there exists a non-negative vector $y$ such that $A^T y = 0$ and $b^T y < 0$.  The slide presents the lemma (Lemma 10.5) and outlines the proof strategy. The proof will involve analyzing a primal linear program (P) which maximizes 0 subject to $Ax \le b$, and its dual linear program (D) which minimizes $b^T y$ subject to $A^T y = 0$ and $y \ge 0$. The equation numbers (10.8) refer to a numbering system within the lecture notes or textbook. The context of linear programming is crucial because Farkas' Lemma is a cornerstone theorem used to establish the duality theory of linear programming and to analyze the feasibility of linear programs.  The lemma provides a powerful tool for determining the solvability of systems of linear inequalities.%
\subsection{Figures}%
\begin{figure}[H]%
\centering%
\begin{subfigure}[b]{0.15\textwidth}%
\centering%
\includegraphics[width=\textwidth]{/Users/ashoksaravanan/Coding/ScribeLec/server/output/MA 421/lectures/2024-10-28-Farkas/figures/1.1.png}%
\caption{A handwritten slide shows the statement of Farkas' Lemma and its proof outline. The lemma states that the system $Ax \le b$ has no solutions if and only if there exists a vector $y$ satisfying three conditions: $A^T y = 0$, $y \ge 0$, and $b^T y < 0$. The proof begins by considering a primal linear program (P) and its dual (D).  The primal program maximizes 0 subject to $Ax \le b$, and the dual minimizes $b^T y$ subject to $A^T y = 0$ and $y \ge 0$.}%
\end{subfigure}%
\end{figure}

%
\newpage%
\section{Slide 2}%
\label{sec:Slide2}%
\subsection{Content}%

\text{Farkas Lemma (Other Versions)} \\
\text{6.4.1 Proposition (Farkas lemma). Let } A \text{ be a real matrix with } m \text{ rows} \\
\text{and } n \text{ columns, and let } b \in \mathbb{R}^m \text{ be a vector. Then exactly one of the following} \\
\text{two possibilities occurs:} \\
\left\{ \begin{array}{l}
(F1) \text{ There exists a vector } x \in \mathbb{R}^n \text{ satisfying } Ax = b \text{ and } x \ge 0. \\
(F2) \text{ There exists a vector } y \in \mathbb{R}^m \text{ such that } y^T A \ge 0^T \text{ and } y^T b < 0.
\end{array} \right. \\
\text{Alternatives}
%
\subsection{Description}%
Slide 2 of 11 presents an alternative version of Farkas' Lemma.  The lemma states that for a given real matrix $A$ with $m$ rows and $n$ columns, and a vector $b \in \mathbb{R}^m$, exactly one of two possibilities holds: (F1) there exists a non-negative vector $x \in \mathbb{R}^n$ such that $Ax = b$, or (F2) there exists a vector $y \in \mathbb{R}^m$ such that $y^T A \ge 0^T$ (all entries of $y^T A$ are non-negative) and $y^T b < 0$.  The notation $0^T$ represents a zero row vector. This version differs from the previous slide's presentation, using equality constraints ($Ax = b$) instead of inequality constraints ($Ax \le b$). The word "Alternatives" is handwritten below, possibly indicating that other equivalent formulations of the lemma exist. The context of linear programming is crucial because this lemma is a fundamental tool for analyzing the feasibility and duality of linear programs.  It provides conditions under which a system of linear equations and non-negativity constraints has a solution.%
\subsection{Figures}%
\begin{figure}[H]%
\centering%
\begin{subfigure}[b]{0.15\textwidth}%
\centering%
\includegraphics[width=\textwidth]{/Users/ashoksaravanan/Coding/ScribeLec/server/output/MA 421/lectures/2024-10-28-Farkas/figures/2.1.png}%
\caption{ }%
\end{subfigure}%
\end{figure}

%
\newpage%
\section{Slide 3}%
\label{sec:Slide3}%
\subsection{Content}%

\text{[MG]} \\
\text{Farkas Lemma (Other Versions)} \\
\text{6.4.1 Proposition (Farkas lemma). Let } A \text{ be a real matrix with } m \text{ rows} \\
\text{and } n \text{ columns, and let } b \in \mathbb{R}^m \text{ be a vector. Then exactly one of the following} \\
\text{two possibilities occurs:} \\
\left\{ \begin{array}{l}
(F1) \text{ There exists a vector } x \in \mathbb{R}^n \text{ satisfying } Ax = b \text{ and } x \ge 0. \\
(F2) \text{ There exists a vector } y \in \mathbb{R}^m \text{ such that } y^T A \ge 0^T \text{ and } y^T b < 0.
\end{array} \right. \\
\text{Alternatives} \\
\underline{\text{Fredholm Alternatives}} \\
AX = b \\
(F1) \quad Ax = b \text{ has a solution} \\
(F2) \quad \text{There is } y \text{ s.t. } y^T A = 0 \text{ and } y^T b \ne 0
%
\subsection{Description}%
Slide 3 of 11 presents alternative formulations of Farkas' Lemma, connecting it to the Fredholm Alternatives. The slide begins by restating Farkas' Lemma (Proposition 6.4.1) in a slightly different form than Slide 2, maintaining the core idea that exactly one of two conditions holds: either a system of equations with non-negativity constraints has a solution (F1), or a related system involving the transpose of the matrix A has a solution (F2).  The handwritten "Alternatives" and underlined "Fredholm Alternatives" suggest a connection between Farkas' Lemma and the Fredholm Alternatives, a concept from functional analysis.  The Fredholm Alternatives are then presented in a simpler form, focusing on the existence of a solution to the system $Ax = b$ (F1) or the existence of a vector $y$ such that $y^T A = 0$ and $y^T b \ne 0$ (F2).  The notation $0^T$ represents a zero row vector. The context of linear programming is crucial because these alternative formulations provide different perspectives on the same fundamental result, offering flexibility in applying the lemma to various linear programming problems. The connection to Fredholm Alternatives highlights the broader mathematical context of Farkas' Lemma.%
\subsection{Figures}%
\begin{figure}[H]%
\centering%
\begin{subfigure}[b]{0.15\textwidth}%
\centering%
\includegraphics[width=\textwidth]{/Users/ashoksaravanan/Coding/ScribeLec/server/output/MA 421/lectures/2024-10-28-Farkas/figures/3.1.png}%
\caption{ }%
\end{subfigure}%
\end{figure}

%
\newpage%
\section{Slide 4}%
\label{sec:Slide4}%
\subsection{Content}%

\text{[MG]} \\
\text{Farkas Lemma (Other Versions)} \\
\text{6.4.3 Proposition (Farkas lemma in three variants). Let } A \text{ be a real} \\
\text{matrix with } m \text{ rows and } n \text{ columns, and let } b \in \mathbb{R}^m \text{ be a vector.} \\
\text{(i) The system } Ax = b \text{ has a nonnegative solution if and only if every} \\
\text{ } y \in \mathbb{R}^m \text{ with } y^T A \ge 0^T \text{ also satisfies } y^T b \ge 0. \\
\text{(ii) The system } Ax \le b \text{ has a nonnegative solution if and only if every} \\
\text{nonnegative } y \in \mathbb{R}^m \text{ with } y^T A \ge 0^T \text{ also satisfies } y^T b \ge 0. \\
\text{(iii) The system } Ax \le b \text{ has a solution if and only if every nonnegative} \\
\text{ } y \in \mathbb{R}^m \text{ with } y^T A = 0^T \text{ also satisfies } y^T b \ge 0. \\
\text{(in fact, (i) } \iff \text{ (ii) } \iff \text{ (iii))} \\
\begin{array}{|c|c|}
\hline
\text{The system} & \text{The system} \\
Ax \le b & Ax = b \\
\hline
\begin{array}{c}
\text{has a solution} \\
x \ge 0 \text{ iff}
\end{array} & \begin{array}{c}
y^T A \ge 0^T \\
\implies y^T b \ge 0
\end{array} \\
\hline
\begin{array}{c}
\text{has a solution} \\
x \in \mathbb{R}^n \text{ iff}
\end{array} & \begin{array}{c}
y \ge 0, y^T A = 0^T \\
\implies y^T b \ge 0
\end{array} \\
\hline
\end{array} \\
\text{Fredholm Alt.}
%
\subsection{Description}%
Slide 4 of 11 presents three variants of Farkas' Lemma, each providing conditions for the solvability of different systems of linear inequalities and equations.  The lemma is stated as Proposition 6.4.3.  The three variants concern: (i) the existence of a non-negative solution to $Ax = b$; (ii) the existence of a non-negative solution to $Ax \le b$; and (iii) the existence of any solution to $Ax \le b$.  Each variant provides a necessary and sufficient condition involving a dual system of inequalities.  The equivalence of the three variants is explicitly stated.  Two tables summarize the conditions for each variant, clearly showing the relationship between the primal system (involving $x$) and the dual system (involving $y$). The notation $0^T$ represents a zero row vector. The handwritten "Fredholm Alt." likely refers to the Fredholm Alternatives, suggesting a connection between Farkas' Lemma and this concept from functional analysis. The context of linear programming is crucial because these variants of Farkas' Lemma are fundamental tools for analyzing the feasibility and duality of linear programs. They provide powerful criteria for determining the existence of solutions to various linear systems arising in optimization problems.%
\subsection{Figures}%
\begin{figure}[H]%
\centering%
\begin{subfigure}[b]{0.15\textwidth}%
\centering%

\includegraphics[width=\textwidth]{/Users/ashoksaravanan/Coding/ScribeLec/server/output/MA 421/lectures/2024-10-28-Farkas/figures/4.1.png}%
\caption{A table summarizing three variants of Farkas' Lemma.}%
\end{subfigure}%
\hspace{0.05\textwidth}%
\begin{subfigure}[b]{0.15\textwidth}%
\centering%

\includegraphics[width=\textwidth]{/Users/ashoksaravanan/Coding/ScribeLec/server/output/MA 421/lectures/2024-10-28-Farkas/figures/4.2.png}%
\caption{Another table summarizing three variants of Farkas' Lemma.}%
\end{subfigure}%
\end{figure}

%
\newpage%
\section{Slide 5}%
\label{sec:Slide5}%
\subsection{Content}%

\text{[MG]} \\
\text{Farkas Lemma (Other Versions)} \\
\text{6.4.3 Proposition (Farkas lemma in three variants). Let } A \text{ be a real} \\
\text{matrix with } m \text{ rows and } n \text{ columns, and let } b \in \mathbb{R}^m \text{ be a vector.} \\
\text{(i) The system } Ax = b \text{ has a nonnegative solution if and only if every} \\
\text{ } y \in \mathbb{R}^m \text{ with } y^T A \ge 0^T \text{ also satisfies } y^T b \ge 0. \\
\text{(ii) The system } Ax \le b \text{ has a nonnegative solution if and only if every} \\
\text{nonnegative } y \in \mathbb{R}^m \text{ with } y^T A \ge 0^T \text{ also satisfies } y^T b \ge 0. \\
\text{(iii) The system } Ax \le b \text{ has a solution if and only if every nonnegative} \\
\text{ } y \in \mathbb{R}^m \text{ with } y^T A = 0^T \text{ also satisfies } y^T b \ge 0. \\
\text{(in fact, (i) } \iff \text{ (ii) } \iff \text{ (iii))} \\
\sim \text{(i) } Ax = b, x \ge 0 \text{ has not solution} \\
\text{iff there is } y \text{ s.t. } y^T A > 0^T, y^T b < 0
%
\subsection{Description}%
Slide 5 of 11 presents three variants of Farkas' Lemma and an additional handwritten statement related to the non-existence of a solution.  Proposition 6.4.3 formally states three equivalent versions of the lemma, each concerning the solvability of a linear system ($Ax = b$, $Ax \le b$) under different conditions on the solution vector $x$ and a dual system involving $y$.  The conditions involve non-negativity constraints and inequalities on $y^T A$ and $y^T b$. The notation $0^T$ represents a zero row vector. The handwritten part below states that the system $Ax = b$ with $x \ge 0$ (non-negative $x$) has no solution if and only if there exists a vector $y$ such that $y^T A > 0^T$ (all entries of $y^T A$ are strictly positive) and $y^T b < 0$. This provides a condition for the *non*-existence of a solution, complementing the existence conditions given in the proposition. The context of linear programming is crucial because these variants and the additional statement provide powerful tools for analyzing the feasibility of linear programs. They offer different perspectives on the same fundamental result, allowing for flexibility in applying the lemma to various linear programming problems.%
\subsection{Figures}%
\begin{figure}[H]%
\centering%
\begin{subfigure}[b]{0.15\textwidth}%
\centering%

\includegraphics[width=\textwidth]{/Users/ashoksaravanan/Coding/ScribeLec/server/output/MA 421/lectures/2024-10-28-Farkas/figures/5.1.png}%
\caption{ }%
\end{subfigure}%
\end{figure}

%
\newpage%
\section{Slide 6}%
\label{sec:Slide6}%
\subsection{Content}%

\text{[MG]} \\
\text{Farkas Lemma (Other Versions)} \\
\text{6.4.3 Proposition (Farkas lemma in three variants). Let } A \text{ be a real} \\
\text{matrix with } m \text{ rows and } n \text{ columns, and let } b \in \mathbb{R}^m \text{ be a vector.} \\
\text{(i) The system } Ax = b \text{ has a nonnegative solution if and only if every} \\
\text{ } y \in \mathbb{R}^m \text{ with } y^T A \ge 0^T \text{ also satisfies } y^T b \ge 0. \\
\text{(ii) The system } Ax \le b \text{ has a nonnegative solution if and only if every} \\
\text{nonnegative } y \in \mathbb{R}^m \text{ with } y^T A \ge 0^T \text{ also satisfies } y^T b \ge 0. \\
\text{(iii) The system } Ax \le b \text{ has a solution if and only if every nonnegative} \\
\text{ } y \in \mathbb{R}^m \text{ with } y^T A = 0^T \text{ also satisfies } y^T b \ge 0. \\
\text{(in fact, (i) } \iff \text{ (ii) } \iff \text{ (iii))} \\
\sim \text{(ii) } Ax \le b, x \ge 0 \text{ has no solution} \\
\text{iff there is } y \ge 0 \text{ s.t. } y^T A \ge 0^T, y^T b < 0
%
\subsection{Description}%
Slide 6 of 11 presents three variants of Farkas' Lemma and a handwritten statement concerning the non-existence of a solution. Proposition 6.4.3 formally states three equivalent versions of the lemma, each concerning the solvability of a linear system ($Ax = b$, $Ax \le b$) under different conditions on the solution vector $x$ and a dual system involving $y$. The conditions involve non-negativity constraints and inequalities on $y^T A$ and $y^T b$. The notation $0^T$ represents a zero row vector. The handwritten part below states that the system $Ax \le b$ with $x \ge 0$ (non-negative $x$) has no solution if and only if there exists a non-negative vector $y$ such that $y^T A \ge 0^T$ and $y^T b < 0$. This provides a condition for the *non*-existence of a solution, complementing the existence conditions given in the proposition.  The use of $\sim$(ii) indicates this is a related statement, likely derived from or closely connected to variant (ii) of the proposition. The context of linear programming is crucial because these variants and the additional statement provide powerful tools for analyzing the feasibility of linear programs. They offer different perspectives on the same fundamental result, allowing for flexibility in applying the lemma to various linear programming problems.%
\subsection{Figures}%
\begin{figure}[H]%
\centering%
\begin{subfigure}[b]{0.15\textwidth}%
\centering%
\includegraphics[width=\textwidth]{/Users/ashoksaravanan/Coding/ScribeLec/server/output/MA 421/lectures/2024-10-28-Farkas/figures/6.1.png}%
\caption{ }%
\end{subfigure}%
\end{figure}

%
\newpage%
\section{Slide 7}%
\label{sec:Slide7}%
\subsection{Content}%

\text{[MG]} \\
\text{Farkas Lemma (Other Versions)} \\
\text{6.4.3 Proposition (Farkas lemma in three variants). Let } A \text{ be a real} \\
\text{matrix with } m \text{ rows and } n \text{ columns, and let } b \in \mathbb{R}^m \text{ be a vector.} \\
\text{(i) The system } Ax = b \text{ has a nonnegative solution if and only if every} \\
\text{ } y \in \mathbb{R}^m \text{ with } y^T A \ge 0^T \text{ also satisfies } y^T b \ge 0. \\
\text{(ii) The system } Ax \le b \text{ has a nonnegative solution if and only if every} \\
\text{nonnegative } y \in \mathbb{R}^m \text{ with } y^T A \ge 0^T \text{ also satisfies } y^T b \ge 0. \\
\text{(iii) The system } Ax \le b \text{ has a solution if and only if every nonnegative} \\
\text{ } y \in \mathbb{R}^m \text{ with } y^T A = 0^T \text{ also satisfies } y^T b \ge 0. \\
\text{(in fact, (i) } \iff \text{ (ii) } \iff \text{ (iii))} \\
\sim \text{(iii) } Ax \le b \text{ has no solution} \\
\text{iff there is } y \text{ s.t. } y^T A = 0^T, y^T b < 0
%
\subsection{Description}%
Slide 7 of 11 continues the discussion of Farkas' Lemma, focusing on alternative formulations. Proposition 6.4.3 presents three equivalent versions of the lemma, each providing conditions for the solvability of linear systems ($Ax = b$, $Ax \le b$) under different constraints on the solution vector $x$ and a dual system involving $y$.  The conditions involve non-negativity constraints and inequalities on $y^T A$ and $y^T b$. The notation $0^T$ represents a zero row vector.  The handwritten section below presents a statement related to the *non*-existence of a solution to the system $Ax \le b$. It states that this system has no solution if and only if there exists a vector $y$ such that $y^T A = 0^T$ and $y^T b < 0$. This condition is a direct consequence of, or closely related to, variant (iii) of the proposition. The symbol $\sim$ indicates a connection to variant (iii). The context of linear programming is crucial because these variants and the additional statement provide powerful tools for analyzing the feasibility of linear programs. They offer different perspectives on the same fundamental result, allowing for flexibility in applying the lemma to various linear programming problems.%
\subsection{Figures}%
\begin{figure}[H]%
\centering%
\begin{subfigure}[b]{0.15\textwidth}%
\centering%
\includegraphics[width=\textwidth]{/Users/ashoksaravanan/Coding/ScribeLec/server/output/MA 421/lectures/2024-10-28-Farkas/figures/7.1.png}%
\caption{ }%
\end{subfigure}%
\end{figure}

%
\newpage%
\section{Slide 8}%
\label{sec:Slide8}%
\subsection{Content}%

\text{[MG]} \\
\text{Farkas Lemma (Other Versions)} \\
\text{6.4.1 Proposition (Farkas lemma). Let } A \text{ be a real matrix with } m \text{ rows} \\
\text{and } n \text{ columns, and let } b \in \mathbb{R}^m \text{ be a vector. Then exactly one of the following} \\
\text{two possibilities occurs:} \\
\begin{cases}
\text{(F1) There exists a vector } x \in \mathbb{R}^n \text{ satisfying } Ax = b \text{ and } x \ge 0. \\
\text{(F2) There exists a vector } y \in \mathbb{R}^m \text{ such that } y^T A \ge 0^T \text{ and } y^T b < 0.
\end{cases} \\
\text{Alternatives}
%
\subsection{Description}%
Slide 8 of 11 presents Proposition 6.4.1, a version of Farkas' Lemma.  The proposition states that for a given real matrix $A$ ($m$ rows, $n$ columns) and vector $b \in \mathbb{R}^m$, exactly one of two possibilities holds: (F1) there exists a non-negative solution $x \in \mathbb{R}^n$ to the system $Ax = b$, or (F2) there exists a vector $y \in \mathbb{R}^m$ such that $y^T A \ge 0^T$ (all entries of $y^T A$ are non-negative) and $y^T b < 0$.  The notation $0^T$ represents a zero row vector.  The handwritten "Alternatives" suggests that other versions or formulations of Farkas' Lemma will follow. The context is linear programming; this lemma is a fundamental tool for analyzing the feasibility and duality of linear programs.  The lemma establishes a strong duality relationship: if a primal problem is infeasible, then its dual problem is unbounded, and vice versa.  This specific formulation focuses on the existence of a non-negative solution to a system of linear equations.%
\subsection{Figures}%
\begin{figure}[H]%
\centering%
\begin{subfigure}[b]{0.15\textwidth}%
\centering%
\includegraphics[width=\textwidth]{/Users/ashoksaravanan/Coding/ScribeLec/server/output/MA 421/lectures/2024-10-28-Farkas/figures/8.1.png}%
\caption{ }%
\end{subfigure}%
\end{figure}

%
\newpage%
\section{Slide 9}%
\label{sec:Slide9}%
\subsection{Content}%

\text{[MG]} \\
\text{Farkas Lemma (Other Versions)} \\
\text{6.4.1 Proposition (Farkas lemma). Let } A \text{ be a real matrix with } m \text{ rows} \\
\text{and } n \text{ columns, and let } b \in \mathbb{R}^m \text{ be a vector. Then exactly one of the following} \\
\text{two possibilities occurs:} \\
\begin{cases}
\text{(F1) There exists a vector } x \in \mathbb{R}^n \text{ satisfying } Ax = b \text{ and } x \ge 0. \\
\text{(F2) There exists a vector } y \in \mathbb{R}^m \text{ such that } y^T A \ge 0^T \text{ and } y^T b < 0.
\end{cases} \\
\text{Alternatives} \\
\text{The plan of the proof is straightforward: We let } z \text{ be the point of } C \text{ nearest} \\
\text{to } b \text{ (in the Euclidean distance), and we check that the vector } y = z - b \text{ is} \\
\text{as required; see the following illustration:} \\
\text{Cone: } \{Ax \ge 0\}
%
\subsection{Description}%
Slide 9 of 11 presents Proposition 6.4.1, a statement of Farkas' Lemma.  It states that for a given matrix $A$ ($m \times n$) and vector $b \in \mathbb{R}^m$, exactly one of two conditions holds: (F1) a non-negative solution $x \in \mathbb{R}^n$ exists for $Ax = b$, or (F2) a vector $y \in \mathbb{R}^m$ exists such that $y^T A \ge 0^T$ and $y^T b < 0$. The handwritten "Alternatives" suggests that other versions of the lemma will be discussed. The slide then outlines the proof plan: find the point $z$ in the cone $\{Ax \ge 0\}$ closest to $b$ (in Euclidean distance), and show that $y = z - b$ satisfies condition (F2). The figure illustrates this geometrically: the cone $\{Ax \ge 0\}$ is generated by vectors $a_1, a_2, a_3$, point $b$ lies outside the cone, $z$ is the closest point in the cone to $b$, and the vector $y = z - b$ is shown.  The context is linear programming; Farkas' Lemma is a fundamental result used to establish duality theorems and analyze the feasibility of linear programs. The geometric interpretation aids in understanding the lemma's core idea.%
\subsection{Figures}%
\begin{figure}[H]%
\centering%
\begin{subfigure}[b]{0.15\textwidth}%
\centering%
\includegraphics[width=\textwidth]{/Users/ashoksaravanan/Coding/ScribeLec/server/output/MA 421/lectures/2024-10-28-Farkas/figures/9.1.png}%
\caption{A geometric illustration of the proof of Farkas' Lemma, showing a cone, a point b outside the cone, and the point z on the cone closest to b. The vector y = z - b is highlighted.}%
\end{subfigure}%
\end{figure}

%
\newpage%
\section{Slide 10}%
\label{sec:Slide10}%
\subsection{Content}%

\text{Farkas Lemma (Other Versions)} \\
\text{[CJ] P.248 \#16.10} \\
\text{Derive the following theorems (with the vector inequality } v > w \text{ meaning, as usual,} \\
v_k > w_k \text{ for all } k \text{) from the result of Problem 16.9.} \\
\text{(i) P. Gordan (1873): The system } Ax < 0 \text{ is unsolvable if and only if the system} \\
yA = 0, y \ge 0, y \ne 0 \text{ is solvable.} \\
\text{(ii) J. Farkas (1902): The system } Ax \le 0, bx > 0 \text{ is unsolvable if and only if the system} \\
yA = b, y \ge 0 \text{ is solvable.} \\
\text{(iii) E. Stiemke (1915): The system } Ax = 0, x > 0 \text{ is unsolvable if and only if the system} \\
yA \ge 0, yA \ne 0 \text{ is solvable.} \\
\text{(iv) J. A. Ville (1938): The system } Ax < 0, x \ge 0 \text{ is unsolvable if and only if the system} \\
yA \ge 0, y \ge 0, y \ne 0 \text{ is solvable.} \\
\text{(v) A. W. Tucker (1956): The system } Ax \ge 0, x \ge 0 \text{ has no solution with } x_k > 0 \text{ if} \\
\text{and only if the system } yA \le 0, y \ge 0 \text{ has a solution with } \sum_{i=1}^m a_{ik} y_i < 0.
%
\subsection{Description}%
Slide 10 of 11 presents a problem (#16.10 from page 248 of a textbook denoted [CJ]) that asks to derive five different theorems related to Farkas' Lemma. Each theorem provides an alternative formulation or a related result concerning the solvability of systems of linear inequalities and equations.  The theorems are attributed to various mathematicians (Gordan, Farkas, Stiemke, Ville, and Tucker) and their respective publication years.  The notation $v > w$ means $v_k > w_k$ for all $k$, indicating element-wise inequality.  Each theorem states that the unsolvability of a certain system of inequalities/equations is equivalent to the solvability of another related system.  The systems involve matrices ($A$), vectors ($x$, $y$, $b$), and inequalities/equations involving non-negativity constraints ($x \ge 0$, $y \ge 0$) and strict inequalities ($x > 0$, $bx > 0$). The context is linear programming; these theorems are variations of Farkas' Lemma and are fundamental tools for analyzing the duality and feasibility of linear programs.  The problem requires proving the equivalence of the solvability conditions stated in each theorem, likely using a previously established version of Farkas' Lemma (as referenced by "Problem 16.9").%
\subsection{Figures}%
\begin{figure}[H]%
\centering%
\begin{subfigure}[b]{0.15\textwidth}%
\centering%

\includegraphics[width=\textwidth]{/Users/ashoksaravanan/Coding/ScribeLec/server/output/MA 421/lectures/2024-10-28-Farkas/figures/10.1.png}%
\caption{ }%
\end{subfigure}%
\end{figure}

%
\newpage%
\section{Slide 11}%
\label{sec:Slide11}%
\subsection{Content}%

\text{[MG]} \\
\text{Summary of Farkas Lemma}
%
\subsection{Description}%
Slide 11 of 11 provides a summary of Farkas' Lemma.  The handwritten "[MG]" likely indicates the author's initials or a source reference. The slide title, "Summary of Farkas Lemma," suggests that this slide concludes a discussion on Farkas' Lemma, a fundamental result in linear programming.  The slide itself is blank, indicating that the summary was likely presented verbally or through additional material not captured in the image.  The context of the course (Linear Programming) is crucial because Farkas' Lemma is a key theorem used to prove duality theorems and analyze the feasibility of linear programs.  The absence of content on the slide implies that the preceding slides contained the detailed statements and proofs of Farkas' Lemma and its various equivalent forms.%
\subsection{Figures}%
\begin{figure}[H]%
\centering%
\begin{subfigure}[b]{0.15\textwidth}%
\centering%
\includegraphics[width=\textwidth]{/Users/ashoksaravanan/Coding/ScribeLec/server/output/MA 421/lectures/2024-10-28-Farkas/figures/11.1.png}%
\caption{ }%
\end{subfigure}%
\end{figure}

%
\newpage%
\end{document}